\chapter[Conclusão]{Conclusão}
\label{cap:conclusao}

Este trabalho abordou a extração de relações entre entidades clínicas em português
sobre o corpus SemClinBr, no espaço de rótulos \texttt{negation\_of},
\texttt{associated\_with} e \texttt{no\_relation}. Como etapa inicial do
desenvolvimento do \textbf{RECLin-PT}, conduziu-se um experimento controlado de
\emph{baselines} para responder a uma pergunta que fundamenta as decisões de
projeto seguintes: o pré-treinamento clínico importa para a extração de relações em
textos médicos em português? A ênfase recaiu, simultaneamente, sobre a resposta a
essa pergunta e sobre a construção de uma cadeia experimental reprodutível, do
texto bruto às métricas finais.

\section{Contribuições}

As principais contribuições são: (i) um \emph{parser} fiel e determinístico do
formato XML do SemClinBr, com verificação de sanidade e correção controlada de
\emph{offsets}; (ii) uma caracterização do corpus que evidencia o desbalanceamento
das relações e a proximidade entre entidades relacionadas, em particular a
quase adjacência dos pares de negação; (iii) um protocolo de particionamento em
nível de documento, estratificado e congelado, que previne vazamento de dados;
(iv) um gerador determinístico de pares candidatos compartilhado por todos os
modelos, com o teto de \emph{recall} imposto pela janela quantificado por partição
e por classe; (v) uma infraestrutura experimental de paridade total, em que os dois
\emph{encoders} comparados diferem apenas no \emph{checkpoint} de pré-treino; e
(vi) uma comparação empírica com teste de significância pareado entre um
\emph{encoder} clínico e um geral sobre a mesma tarefa.

Quanto aos resultados, não se observou vantagem \emph{consistente} do
pré-treinamento clínico. Na semente de referência ($42$), a diferença no F1 de
\texttt{negation\_of} entre BioBERTpt e BERTimbau é de $0{,}017$ a favor do geral e
não é estatisticamente significativa na métrica-alvo (\emph{bootstrap} pareado
IC95\% $[-0{,}053;\,+0{,}018]$, $p=0{,}345$); na semente $43$, o mesmo
\emph{bootstrap} sobre essa métrica inverte o sinal e passa a favorecer o BioBERTpt
clínico, agora de forma significativa (IC95\% $[+0{,}016;\,+0{,}090]$,
$p=0{,}0054$). Nem o veredito estatístico nem a \emph{direção} da diferença
\textbf{na classe-alvo}, portanto, são estáveis entre sementes, pois a variação de
inicialização de um mesmo modelo ($0{,}078$ de amplitude no BioBERTpt) supera com
folga a diferença entre os dois modelos. Essa instabilidade é da métrica-alvo e
\textbf{não se estende ao padrão global de acertos}. O teste de McNemar, que incide
sobre as três classes, rejeita a hipótese nula nas duas sementes e aponta nas duas
para o \emph{encoder} geral ($540$ contra $695$ discordâncias na semente $42$; $442$
contra $624$ na $43$). O que é estável, além dessa direção agregada, é a assimetria
de dispersão. O \emph{encoder} geral acompanha o clínico em macro-F1 médio
($0{,}695$ contra $0{,}694$) com cerca de $60\%$ dos parâmetros e com desempenho
cerca de dez vezes mais estável entre reinicializações na métrica-alvo.

\section{Trabalhos futuros}

Entregar um estudo de \emph{baselines} e deixar a especialização do modelo para a fase
seguinte é uma decisão de sequenciamento, e não uma lacuna. A escolha do \emph{encoder}
é pré-requisito de qualquer aprimoramento posterior. Fixada a base por expectativa
(e a expectativa de partida, registrada na Seção~\ref{sec:motivacao}, era a de
vantagem do pré-treino clínico), todo ganho ou perda medido depois ficaria
confundido com o efeito do \emph{checkpoint}, sem meio de separar um do outro. E o
resultado obtido é, ele próprio, um achado de consequência prática, e não um resultado
nulo vazio. Sob teste pareado e sob duas sementes de treino, o \emph{encoder} clínico
não leva vantagem \emph{confiável} nesta tarefa (a diferença medida troca de sinal
entre sementes), o que libera a construção do RECLin-PT sobre um modelo
sensivelmente menor. Isso não equivale a custo zero, pois, tomada a média das duas
sementes, o \emph{encoder} clínico fica $0{,}017$ à frente na métrica-alvo
($0{,}715$ contra $0{,}698$), e é esse o custo esperado da escolha. Ele é, contudo,
menor que a própria amplitude entre sementes do BioBERTpt ($0{,}078$), o que o
torna mal estimado por duas execuções. Como registra o
Capítulo~\ref{cap:relacionados}, não localizamos avaliação controlada equivalente para
o português, de modo que esse achado não era obtenível por consulta à literatura.

Esses resultados orientam diretamente a fase seguinte. O RECLin-PT não é um trabalho a
iniciar. Os \emph{baselines} aqui reportados são a sua primeira fase, e o que
constitui trabalho futuro imediato é a especialização do modelo sobre a base
escolhida. Como o BERTimbau geral acompanha o clínico em macro-F1 médio, exige cerca de $60\%$
dos parâmetros e varia cerca de dez vezes menos entre sementes na métrica-alvo, ele se
apresenta como a base mais atraente por eficiência e previsibilidade, ressalvado
que a semente $43$ mostra o clínico à frente na negação, de modo que a escolha se
apoia em custo e estabilidade, e não em superioridade média demonstrada. Esse refinamento terá como foco o
aprimoramento da classe \texttt{negation\_of}, a mais crítica para a correta
interpretação clínica.

Cabe explicitar que essa prioridade não recai sobre a classe de pior desempenho. A
Seção~\ref{sec:resultados} mostra que \texttt{associated\_with} é a mais difícil das
três (F1 $\approx 0{,}42$--$0{,}43$, contra $\approx 0{,}68$--$0{,}69$ de
\texttt{negation\_of}). O critério é de criticidade clínica, e não de facilidade. Uma
negação lida como afirmação inverte o sentido de um achado no prontuário, ao passo que
um \texttt{associated\_with} não recuperado empobrece a extração sem alterar o sentido
do que foi extraído. O teto de \emph{recall} de \texttt{associated\_with} permanece,
ainda assim, entre as continuidades previstas a seguir.

Como continuidade, prevê-se ainda: estimar o desempenho dos \emph{baselines} sob um
número maior de sementes de treino, com agregação formal, para obter um intervalo
de confiança da diferença mais estreito; enriquecer a representação de entrada com
o \textbf{tipo semântico} das entidades, hoje ausente dos marcadores
(Seção~\ref{sec:markers_limitacao}), medindo se o ganho difere entre o
\emph{encoder} clínico e o geral; reavaliar a janela \texttt{max\_gap} da
geração de pares candidatos, elevando o teto de \emph{recall} de
\texttt{associated\_with} sem inviabilizar o custo de treino; investigar o
tratamento explícito do escopo da negação e da direção das relações; estender o
espaço de relações para além das três classes atuais; avaliar modelos maiores e
\emph{encoders} multilíngues \cite{conneau_2020_xlmr}; e integrar uma etapa de
reconhecimento de entidades nomeadas automática, medindo a extração de relações em
um cenário fim a fim.

O trabalho encontra-se, portanto, em andamento. Os \emph{baselines} aqui reportados
estabelecem o ponto de partida do RECLin-PT, e não um encerramento.
